\section*{Bab \ref{ch:b3:homogeneous-catalysis} --- Katalisis Homogen dalam Industri}

\begin{solution}{exo:b3:homogeneous-catalysis:1}
$\ce{RhCl(PPh3)3}$: 16, Rh(I) $d^8$. $\ce{RhCl(PPh3)2}$: 14, Rh(I). $\ce{RhClH2(PPh3)2}$: 16, Rh(III) $d^6$.
Kompleks alkena: 18, Rh(III). Alkil hidrida: 16, Rh(III). Eliminasi reduktif
kembali ke 14, Rh(I).
\end{solution}

\begin{solution}{exo:b3:homogeneous-catalysis:2}
Butanal, $\ce{CH3CH2CH2CHO}$ (linear), dan 2-metilpropanal, $\ce{(CH3)2CHCHO}$ (bercabang).
\end{solution}

\begin{solution}{exo:b3:homogeneous-catalysis:3}
Dietil siklopent-3-ena-1,1-dikarboksilat, sebuah cincin beranggota lima, dengan lepasnya
etena: kedua gugus $\ce{CH2}$ ujung lepas sebagai $\ce{C2H4}$ dan kedua karbon CH di bagian dalam
tersambung.
\end{solution}

\begin{solution}{exo:b3:homogeneous-catalysis:4}
Metil (\textit{E})-3-fenilprop-2-enoat (metil sinamat): aril teradisi pada
karbon ujung, dan eliminasi $\beta$-hidrida, secara sin, dari konformer yang lebih stabil
memberikan alkena \textit{E}.
\end{solution}

\begin{solution}{exo:b3:homogeneous-catalysis:5}
Tahap adisi oksidatifnya lambat, sehingga keadaan istirahat katalisnya adalah $\ce{[Rh(CO)2I2]-}$ dan
$v = k[\mathrm{Rh}]_{\mathrm{tot}}[\ce{CH3I}]$. Metanol hanya meregenerasi metil iodida melalui reaksi cepat dengan HI;
iodida total yang dimasukkan menentukan $[\ce{CH3I}]$, sehingga laju tidak bergantung pada konsentrasi
metanol sampai metanol yang tersisa terlalu sedikit untuk mengubah HI kembali menjadi
metil iodida.
\end{solution}

\begin{solution}{exo:b3:homogeneous-catalysis:6}
Asam fenilboronat dengan 4-bromotoluena, atau asam 4-metilfenilboronat dengan
bromobenzena. Keduanya berhasil; yang pertama memakai asam boronat yang paling murah dan aril bromida
yang stabil dan umum.
\end{solution}

\begin{solution}{exo:b3:homogeneous-catalysis:7}
Bersimetri $C_2$: isotaktik (kedua situs memilih muka yang sama). Bersimetri $C_s$:
sindiotaktik (situs-situs yang saling bayangan cermin berselang-seling). $\ce{Cp2ZrCl2}$ tanpa jembatan: ataktik (situs
akiral).
\end{solution}

\begin{solution}{exo:b3:homogeneous-catalysis:8}
$\mathrm{TON} = 0.98/0.0010 = 980$; rata-rata $\mathrm{TOF} = 980/\qty{2.0}{h} = \qty{490}{h^{-1}}$.
\end{solution}

\begin{solution}{exo:b3:homogeneous-catalysis:9}
Satuan ulangnya adalah $\ce{-[CH=CH-C5H8]-}$ dengan kedua gugus CH pada karbon 1 dan 3 sebuah
cincin siklopentana (poli(1,3-siklopentilenavinilena)). Pelepasan tegangan
cincin bisiklik membuat $\Delta_rH^\circ$ sangat negatif, yang mengalahkan hilangnya
entropi translasi; tidak ada gas yang terbentuk, sehingga gaya dorongnya bersifat
entalpik.
\end{solution}

\begin{solution}{exo:b3:homogeneous-catalysis:10}
$\bar n_n = 1/(1 - p) = 2000$; $M_w/M_n = 1 + p = 1.9995 \approx 2.0$. \omterm{def:b3:homogeneous-catalysis:ziegler-natta}{Katalis Ziegler--Natta} memiliki
beberapa jenis situs, masing-masing dengan $p$-nya sendiri; jumlah beberapa distribusi Schulz--Flory
dengan rata-rata yang berbeda lebih lebar (dispersitas sering 4 sampai 8).
\end{solution}

\begin{solution}{exo:b3:homogeneous-catalysis:11}
Pada tekanan rendah $K_1K_2[\ce{H2}] \ll [\mathrm L] + K_1$: $v \approx kK_1K_2[\mathrm{Rh}]_{\mathrm{tot}}[\text{alkena}][\ce{H2}]/([\mathrm L] + K_1)$,
berorde satu terhadap $\ce{H2}$. Pada tekanan tinggi suku $K_1K_2[\ce{H2}]$ mendominasi: $v \to k[\mathrm{Rh}]_{\mathrm{tot}}[\text{alkena}]$,
tidak bergantung pada $\ce{H2}$. $[\mathrm L]$ hanya muncul di penyebut: fosfina yang ditambahkan menurunkan $v$
dengan mendorong rodium kembali ke prakatalis yang jenuh.
\end{solution}

\begin{solution}{exo:b3:homogeneous-catalysis:12}
Insersi yang menempatkan rodium pada karbon dalam membangun alkil sekunder
di samping logam yang sesak; fosfina yang meruah (dan kelebihannya, yang menjaga dua fosfina
tetap pada logam) membuat keadaan transisi itu lebih buruk daripada keadaan transisi yang memberikan alkil
primer, yang menghasilkan aldehida linear. Butanal linear itulah yang diubah melalui
kondensasi aldol dan hidrogenasi menjadi alkohol $C_8$ bercabang yang dipakai dalam
pemlastis; 2-metilpropanal adalah produk samping yang nilainya lebih rendah.
\end{solution}

\begin{solution}{pb:b3:homogeneous-catalysis:1}
\textbf{1.} Ir(I), $d^8$: $9 + 4 + 2 + 1 = 16$ elektron (penghitungan netral, dengan memperhitungkan muatannya).

\textbf{2.} $\ce{[CH3Ir(CO)2I3]-}$, 18 elektron, Ir(III).

\textbf{3.} $\ce{[CH3Ir(CO)3I2]}$, 18 elektron, Ir(III).

\textbf{4.} $\ce{[CH3COIr(CO)2I2]}$, 16 elektron, Ir(III).

\textbf{5.} $\ce{[CH3COIr(CO)2I3]-}$ (18) mengeliminasi $\ce{CH3COI}$, sehingga $\ce{[Ir(CO)2I2]-}$ teregenerasi.

\textbf{6.} $\ce{CH3COI + H2O -> CH3COOH + HI}$; $\ce{CH3OH + HI -> CH3I + H2O}$.

\textbf{7.} $\ce{CH3OH + CO -> CH3COOH}$.

\textbf{8.} Dengan iridium, insersi migrasinya lambat; adisi oksidatif
metil iodida, yang merupakan tahap penentu laju dengan rodium, jauh lebih cepat pada iridium yang lebih
kaya elektron.

\textbf{9.} \omterm{def:b3:surfaces-catalysis:catalyst-design}{Promotor} itu melepaskan satu iodida dari $\ce{[CH3Ir(CO)2I3]-}$, sehingga terbentuk trikarbonil netral, yang
logamnya lebih sedikit berdonasi balik dan gugus metilnya bermigrasi lebih cepat ke CO yang lebih
elektrofilik.

\textbf{10.} Ketergantungan terbalik pada konsentrasi iodida bebas: iodida mendorong
prakesetimbangan kembali ke kompleks anionik yang lambat.

\textbf{11.} Metanol baru masuk sesudah tahap penentu laju, dan diubah menjadi
metil iodida melalui reaksi cepat dengan HI.

\textbf{12.} $M(\ce{CH3COOH}) = \qty{60.1}{g/mol}$: $n = \qty{5.0e11}{g}/\qty{60.1}{g/mol} = \qty{8.3e9}{mol}$ per tahun.

\textbf{13.} $\qty{8.32e9}{mol}/\qty{8000}{h} = \qty{1.04e6}{mol/h}$.

\textbf{14.} $\qty{2.0e5}{L} \times \qty{5.0e-3}{mol/L} = \qty{1.0e3}{mol}$ iridium, $\qty{1.0e3}{mol} \times \qty{192.2}{g/mol} \approx \qty{190}{kg}$.

\textbf{15.} $\mathrm{TOF} = \qty{1.04e6}{mol/h}/\qty{1.0e3}{mol} = \qty{1.0e3}{h^{-1}}$, sekitar \qty{0.29}{s^{-1}}.

\textbf{16.} $\mathrm{TON} = \qty{8.3e9}{mol}/\qty{1.0e3}{mol} = \num{8.3e6}$ per tahun.

\textbf{17.} $\qty{5.0e8}{kg}/(\qty{8000}{h} \times \qty{200}{m^3}) \approx \qty{310}{kg.m^{-3}.h^{-1}}$.

\textbf{18.} $\ce{CO + H2O -> CO2 + H2}$.

\textbf{19.} $1/0.94 = \qty{1.064}{mol}$ CO per mol asam asetat.

\textbf{20.} $1.064 - 1 = \qty{0.064}{mol}$ \ce{CO2} per mol asam asetat.

\textbf{21.} Satu ton adalah $\qty{1.0e6}{g}/\qty{60.1}{g/mol} = \qty{1.66e4}{mol}$; $\qty{1.66e4}{mol} \times 0.0638 \times \qty{44.0}{g/mol} \approx \qty{47}{kg}$ \ce{CO2}.

\textbf{22.} Air adalah pereaksi dalam reaksi geser: air yang lebih sedikit memperlambatnya dan menghemat energi
untuk mengeringkan asam. Air tetap diperlukan untuk menghidrolisis asetil iodida dan untuk menjaga
katalis tetap aktif dan terlarut; dengan rodium, air yang sedikit mengendapkan rodium
iodida.

\textbf{23.} Hidrogen: gas itu menumpuk dalam fase gas, menurunkan tekanan parsial CO
(sehingga gas harus dibuang, dan CO ikut hilang), dan dapat menghidrogenasi zat-zat antara menjadi
produk samping seperti asam propionat.

\textbf{24.} Katalis iridium berputar sekitar $\qty{1.0e3}{h^{-1}}$: setiap atom iridium membuat
seribu molekul asam asetat per jam, sekitar delapan juta per tahun.
\end{solution}
